\documentclass[10pt]{beamer}

\usepackage[T1]{fontenc}
\usepackage{FiraSans} 


%
% Choose how your presentation looks.
%
% For more themes, color themes and font themes, see:
% http://deic.uab.es/~iblanes/beamer_gallery/index_by_theme.html
%
\mode<presentation>
{
  \usetheme[progressbar=foot,numbering=fraction,background=light]{metropolis}      % or try Darmstadt, Madrid, Warsaw, ...
  \usecolortheme{default} % or try albatross, beaver, crane, ...
  \usefonttheme{default}  % or try serif, structurebold, ...
  \setbeamertemplate{navigation symbols}{}
  \setbeamertemplate{caption}[numbered]
  %\setbeamertemplate{frame footer}{My custom footer}
} 

\usepackage[ngerman]{babel}
\usepackage[utf8x]{inputenc}


%--------------------------Editor mode.

\usepackage
[citestyle=authoryear,
%style=historian, 		%Loads the Historian files
sorting=nty,	  		%Sorts bibliography by year, name, title
autocite=footnote, 		%Autocite command generates footnotes
autolang=hyphen, 			%Allows hyphenation rules for foreign languages to apply to individual entries.						%(The other language rules should all be American)
mincrossrefs=1, 		%Includes all x-ref’ed entries in the bibliography
%usetranslator=true, 	%Translator’s name may be substituted for						%author or editor, if the latter are blank
%printseries, 			%Options provided by Historian, see below
backend=biber]
{biblatex}

\DeclareFieldFormat{postnote}{#1}
\DeclareFieldFormat{multipostnote}{#1}
\DeclareAutoCiteCommand{footnote}[f]{\footcite}{\footcites}

\bibliography{literature}
%----------------------------------------

%--------------------------Editor mode.


%\usepackage[default]{raleway}
\usepackage{fontawesome}


\usepackage{hyperref}
%\usepackage{enumitem} % produces fatal error with this template

\usepackage{xcolor}
\definecolor{customcolor}{HTML}{616AC5}
\definecolor{alert}{HTML}{CD5C5C}
\definecolor{w3schools}{HTML}{4CAF50}
\definecolor{subbox}{gray}{0.60}
\definecolor{codecolor}{HTML}{FFC300}



 
%Information to be included in the title page:
\title[XPath \& XSLT] %optional
{LV Textkodierung und Textanalyse 
mit TEI, UE}
\subtitle{Einheit 4: XPath und XSLT-Basisstylesheets}
\institute{Zentrum für Informationsmodellierung, Graz}
\author[SL]{Sarah Lang}
\date[2002119] % (optional)
{Universität Wien, Jänner 2021}

\logo{%\includegraphics[height=1cm]{unipassau.png}
%\includegraphics[height=1cm]{univie-logo.png}
%\includegraphics[height=2cm]{zim.png}
}



\newcommand{\punkti}{~\lbrack\dots\rbrack~}

\renewenvironment{quote}
               {\list{\faQuoteLeft\phantom{ }}{\rightmargin\leftmargin}%
                \item\relax\footnotesize\ignorespaces}
               {\unskip\unskip\phantom{xx}\faQuoteRight\endlist}

\newcommand{\bgupper}[3]{\colorbox{#1}{\color{#2}\huge\bfseries\MakeUppercase{#3}}}
\newcommand{\bg}[3]{\colorbox{#1}{\bfseries\color{#2}#3}}

\newcommand{\mycommand}[2]{{\ttfamily\detokenize{#1}}~\dotfill{}~{\footnotesize #2}\\}

\newcommand{\sep}{{\scriptsize~\faCircle{ }~}}

\newcommand{\red}[1]{\bg{alert}{white}{#1}\\}
\newcommand{\green}[1]{\bg{w3schools}{white}{#1}\\}




%----------------------------------------------------------------------------------------------------------------


\usepackage{xcolor}
\definecolor{customcolor}{HTML}{616AC5}
\definecolor{alert}{HTML}{CD5C5C}
\definecolor{w3schools}{HTML}{4CAF50}
\definecolor{subbox}{gray}{0.60}
\definecolor{codecolor}{HTML}{FFC300}


%--------------------------------------------------------------------------------
\usepackage{tcolorbox}

\tcbuselibrary{most,listingsutf8,minted}

\tcbset{tcbox width=auto,left=1mm,top=1mm,bottom=1mm,
right=1mm,boxsep=1mm,middle=1pt}

\newenvironment{mycolorbox}[2]{%
\begin{tcolorbox}[grow to left by=-1em,grow to right by=-1em,capture=minipage,fonttitle=\large\bfseries, enhanced jigsaw,boxsep=1mm,colback=#1!30!white,on line,tcbox width=auto, toptitle=0mm,colframe=#1,opacityback=0.7,nobeforeafter,title=#2]\scriptsize%
}{\end{tcolorbox}\\[0.2em]}

\newenvironment{subbox}[2]{%
\begin{tcolorbox}[capture=minipage,fonttitle=\normalsize\bfseries, enhanced jigsaw,boxsep=1mm,colback=#1!30!white,on line,tcbox width=auto,left=0.3em,top=1mm, toptitle=0mm,colframe=#1,opacityback=0.7,nobeforeafter,title=#2]\scriptsize %
}{\normalsize\end{tcolorbox}\vspace{0.1em}}

\newenvironment{multibox}[1]{%
\begin{tcbraster}[raster columns=#1,raster equal height,nobeforeafter,raster column skip=1em,raster left skip=1em,raster right skip=1em]}{\end{tcbraster}}


\newenvironment{mycodebox}[2]{%
\begin{tcolorbox}[grow to left by=-1em,grow to right by=-1em,capture=minipage,fonttitle=\large\bfseries, enhanced jigsaw,boxsep=1mm,colback=#1!30!white,on line,tcbox width=auto, toptitle=0mm,colframe=#1,opacityback=0.7,nobeforeafter,title=#2]%
}{\end{tcolorbox}\\[0.2em]}

\newtcolorbox{mybox}[2][]{colback=codecolor!10!white,coltitle=red!70!black,
title={#2},fonttitle=\bfseries,#1}


%-------------------------------

\newtcblisting{mypy}[1]{colback=codecolor!5,colframe=codecolor!80!black,listing only, 
minted options={numbers=left, style=tcblatex,fontsize=\scriptsize,breaklines,autogobble,linenos,numbersep=3mm},
left=5mm,enhanced,
title=#1, fonttitle=\bfseries,
listing engine=minted,minted language=python}

\newtcblisting{myxml}[1]{colback=codecolor!5,colframe=codecolor!80!black,listing only, 
minted options={numbers=left, style=tcblatex,fontsize=\tiny,breaklines,autogobble,linenos,numbersep=3mm},
left=5mm,enhanced,
title=#1, fonttitle=\bfseries,
listing engine=minted,minted language=xml}

\newtcblisting{mybiggerxml}[1]{colback=codecolor!5,colframe=codecolor!80!black,listing only, 
minted options={numbers=left, style=tcblatex,fontsize=\footnotesize,breaklines,autogobble,linenos,numbersep=3mm},
left=5mm,enhanced,
title=#1, fonttitle=\bfseries,
listing engine=minted,minted language=xml}

\newtcblisting{myhtml}[1]{colback=codecolor!5,colframe=codecolor!80!black,listing only, 
minted options={numbers=left, style=tcblatex,fontsize=\tiny,breaklines,autogobble,linenos,numbersep=3mm},
left=5mm,enhanced,
title=#1, fonttitle=\bfseries,
listing engine=minted,minted language=html}

\newtcblisting{mycss}[1]{colback=codecolor!5,colframe=codecolor!80!black,listing only, 
minted options={numbers=left, style=tcblatex,fontsize=\tiny,breaklines,autogobble,linenos,numbersep=3mm},
left=5mm,enhanced,
title=#1, fonttitle=\bfseries,
listing engine=minted,minted language=css}


\newtcblisting{myjs}[1]{colback=codecolor!5,colframe=codecolor!80!black,listing only, 
minted options={numbers=left, style=tcblatex,fontsize=\tiny,breaklines,autogobble,linenos,numbersep=3mm},
left=5mm,enhanced,
title=#1, fonttitle=\bfseries,
listing engine=minted,minted language=js}
 
\begin{document}
 
\frame{\titlepage}

\begin{frame}{Übersicht}%\scriptsize
\tableofcontents
\end{frame}
 
\section{Präliminaria}

\begin{frame}{Lehr- \& Lernziele}
\subsection{Lehr- \& Lernziele}
\begin{enumerate}
    \item Kennenlernen von XPath und XSLT
    \item Im Zuge von Übungen sowohl XPath als auch XSLT mindestens 1x selbst verwendet haben 
    \item Überblick, was XPath und XSLT an Mehrwert bieten
    \item Grundsätzliches Verständnis der Funktionsweise von XSLT und XPath
    \item \alert{In diesem Rahmen \emph{nicht} möglich}: Aneignung fundierter Kenntnis und Beherrschung von XSLT und XPath 
\end{enumerate}

\metroset{block=fill}
\begin{alertblock}{Hausübung}
\footnotesize
Eigenständiges Umsetzen (bzw. Modifizieren) einer simplen Transformation in XSLT, im Zuge dessen auch XPath zur Anwendung kommt. Ein Basisstylesheet wird zur Verfügung gestellt, auf dem aufgebaut werden kann.
\end{alertblock}
\end{frame}

\begin{frame}{Vorstellungsrunde}
\subsection{Vorstellungsrunde}
  \begin{columns}[T,onlytextwidth]
    \column{0.35\textwidth}
      \begin{alertblock}{Leitfragen}
      \footnotesize
        \begin{enumerate}
            \item Name, Studienrichtung 
            \item Warum Interesse an Veranstaltung?
            \item Erfahrungen mit Computer und/oder Digital Humanities \& aktueller Kenntnisstand
        \end{enumerate}
      \end{alertblock}

    \column{0.6\textwidth}
      \metroset{block=fill}
      \begin{block}{Sarah Lang}
      \scriptsize
      \begin{itemize}
          \item urspr. aus Aalen (DE/BW)
          \item Studium: Latein \& Französisch \& Geschichte (Lehramt), Bachelor in Archäologie, Master in Religionswissenschaft, KFU Graz
          \item Arbeit am Institut für Religionswissenschaft: u.a. Monster, Horrorfilm - später Alchemie
          \item DH-Zertifikat - Arbeit am Zentrum für Informationsmodellierung (ZIM-ACDH) Graz
          \item Moralische Wochenschriften $\to$ \protect\url{gams.uni-graz.at/mws}
          \item Repositorium zu antiken Fabeln $\to$ \protect\url{gams.uni-graz.at/graf}
          \item DH-Dissertation: Wissensrepräsentation und automatisierte Annotation alchemischer Symbolsprache (Bsp: Korpus des Chymikers Michael Maier, 1568--1622)
      \end{itemize}
      \end{block}
  \end{columns}
\end{frame}



%\section{Unser Thema: Annotation \& Erschließung}
%

\section{Was ist Annotation?}
\begin{frame}[fragile,allowframebreaks]{Annotation und Erschließung}

\bg{alert}{white}{= Annotation mit dem Ziel der Erschließung}\\
irgendwelche oder alle (?) (Meta-)Daten sammeln vs. gezielt spezifische Daten verfügbarmachen \\

\bg{w3schools}{white}{Was ist Annotation?}\\
Zusatzinfo \sep Metadaten \sep Bedeutung kodieren \sep 
Strukturmerkmale, Wissen und Erkenntnisse (implizite Strukturen) für die Maschine explizit machen für Verarbeitung und Nachnutzung. 

\textbf{Im Druckwesen:} inhaltliche Korrektur und typographische Anweisungen (Druckvorbereitung) \sep
\textbf{Computer:} Einbettung von Markup in elektronische Dokumente (log. Struktur beschreiben) \\

\bg{w3schools}{white}{Visual vs. Generic Markup} = Text mit Auszeichnung (Aussage über Bedeutung: Semantik, Struktur, Beziehungen -- William W. Tunnicliffe (1967) \\

\bg{w3schools}{white}{Visual} = presentational \& procedural \\
\bg{w3schools}{white}{Generic} deskriptiv 
\begin{verbatim}
<name type="person">J. W. v. Goethe</name>
\end{verbatim}
$\to$ explizite Marker zur Unterscheidung für die Weiterverarbeitung
\smallskip

Bennenen, charakterisieren, annotieren (auf formalisierte Art und Weise)
\smallskip

\bg{w3schools}{white}{Markup = Annotation = Kodierung }
\end{frame}



\begin{frame}[allowframebreaks]{Stufen der Annotation}

\bg{alert}{white}{Basisannotation}
Formale Kriterien der Textgestalt, z.B. Überschriften, Paragraphen. Hinzufügen von Metadaten (bibliographisch, z.B.).

\bg{alert}{white}{Tiefenerschließung}
Übergang formal-inhaltlich:
Personen, Orte, \emph{Named Entities}, GND -- Zurückgreifen auf Standards und Normdateien zur Vernetzung vieler Projekte.

\bg{alert}{white}{Übergeordnete Forschungsfrage}
Redeanteile der Monster? Wie wird charakterisiert? Kann man Spannung / Grusel `messen'? Gibt es Signalwörter? Visualisierung der Ergebnisse

\end{frame}

%-----------------------------------------------------
\begin{frame}[allowframebreaks]{Markup / Annotation}
\bg{alert}{white}{Annotation = `Auszeichnung' = Mark-up}\\

\begin{quote}
Eine Auszeichnungssprache (\emph{markup language}, abgekürzt ML) ist eine maschinenlesbare Sprache für die Gliederung und Formatierung von Texten und anderen Daten. Der bekannteste Vertreter ist die \emph{Hypertext Markup Language} (HTML), die Kernsprache des World Wide Webs.

Mit Auszeichnungssprachen werden Eigenschaften, Zugehörigkeiten und Darstellungsformen von Abschnitten eines Textes (Zeichen, Wörtern, Absätzen usw. -- ``Elementen'') oder einer Datenmenge beschrieben. Dies geschieht in der Regel, indem sie mit \emph{Tags} markiert werden.\punkti mit der \emph{Standard Generalized Markup Language} (SGML) empfohlene ``Trennung von Struktur und Darstellung''. (\href{https://de.wikipedia.org/wiki/Auszeichnungssprache}{Wikipedia})
\end{quote}

\bg{alert}{white}{maschinenlesbare `Auszeichnung'}\\
\bg{w3schools}{white}{SMGL}~ \bg{w3schools}{white}{HTML}~ \bg{w3schools}{white}{XML}~ \bg{w3schools}{white}{\dots} \\
\bg{alert}{white}{darstellend vs. beschreibend / semantisch:} \\


\begin{itemize}
    \item z.B. `Schriftgröße 14' vs. `Überschrift' (=Typus). (explizit vs. implizit)
    \item Textformatierung vs. Textbedeutung
    \item procedural, representative, descriptive / conceptual (sematisch)
    \item WYSIWYG-Textverarbeitung  vs. WYSIWYM
    \item Vorteil, z.B. bei langen Dokumenten: Aussehen des Elements 1x umdefinieren
    \item verschiedene Repräsentationen eines Dokuments: beschreibend, dargestellt 
    \item Webbrowser `rendern' HTML: in 2 Sichten verfügbar (als Text und dargestellt)
\end{itemize}
\end{frame}
 
 \begin{frame}[allowframebreaks]{Praktische Tipps}
 Tools zum `Switchen' zwischen Markups/Markdown: 
 z.B.~ \href{https://pandoc.org/}{Pandoc} oder \href{http://oxgarage.tei-c.org/#}{OxGarage} (TEI-Fokus). \\
 \footnotesize
 Download-Möglichkeiten überprüfen -- evtl. gibt es das gewünschte Dokument bereits in Basis-Annotation? Hier allerdings potentiellen Korrektur-Aufwand berücksichtigen -- falls Anpassungsbedarf: oft ist man mit einer .txt-Datei von Hand schneller. (Download oder per XSL `herstellen')
 


Binäre Dokumentenformate  wie \emph{.doc, .pdf, .dvi} (\TeX{}-Ausgabeformat) $\neq$ Auszeichnungssprachen.

Ziel: \textbf{Implizites explizieren}.
\end{frame}


%-------------------------------


\section{Annotation im Alltag}
%---------------------------------------------
\begin{frame}[allowframebreaks]{\href{https://de.wikipedia.org/wiki/Standard_Generalized_Markup_Language}{SGML}}
\bg{alert}{white}{Trennung von Inhalt und Darstellung}\\
\bg{alert}{white}{Wie XML: Metasprache} ~
HTML und XML (Nachfolger von SGML genannt $\to$  SGML-basierend) \sep Metasprache zur Definition von Auszeichnungssprachen für Dokumente \sep \emph{Genormte Verallgemeinerte Auszeichnungssprache} (en: \textbf{Standard Generalized Markup Language}

\bgupper{w3schools}{black}{.sgml}\\
\end{frame}





%-----------------------------------------------------
\begin{frame}[fragile,allowframebreaks]{\href{https://de.wikipedia.org/wiki/Rich_Text_Format}{RTF}}
\bg{alert}{white}{Rich Text Format}~
Microsoft 1987 \sep Austauschformat zwischen Textverarbeitungsprogrammen (versch. Hersteller und Betriebssysteme). \sep enthält im Gegensatz zu `plain text' Markup zur Textformatierung\\

\bgupper{w3schools}{black}{.rtf}\\
\footnotesize
\begin{verbatim}
{\rtf1
Guten Tag!
\line
{\i Dies} ist \b{\i ein
\i0 formatierter \b0Text}.
\par
\b Das \b0Ende.
} 
\end{verbatim}\normalsize

\end{frame}

%-----------------------------------------------------
\begin{frame}[fragile,allowframebreaks]{\href{https://de.wikipedia.org/wiki/JavaScript_Object_Notation}{JSON}}

\bg{alert}{white}{JavaScript Object Notation}\\
\bgupper{w3schools}{black}{.json}\\
ausgesprochen wie `Jason' \sep kompaktes Datenformat \sep lesbar \sep key-value-Paare \sep Verschachtelung \\

\bg{alert}{white}{Alternative zu XML? Unterschiede}\\
XML = Struktur-beschreibend, JSON = Syntax-Konvention (nicht deklarativ) \sep JSON: Definition von Instanzen strukturierter Daten \sep JSON = sehr flexibel \sep `lightweight': geriner Overhead, für wenige Daten in strikter Key-Value-Struktur lesbarer \sep valides JavaScript, per \verb|eval()|-Funktion direkt als JS-Objekt umsetzbar

\bg{alert}{white}{Fazit}~ JSON im Vorteil, wo simple key-value-Paare (\textbf{`Einfachheit'}). XML hat und erlaubt mehr \textbf{Komplexität}.
\bg{w3schools}{white}{XML = Auszeichnungssprache}~\\
\bg{w3schools}{white}{JSON = Datenaustauschformat}~

\footnotesize
\begin{verbatim}
{
  "Herausgeber": "Xema",
  "Nummer": "1234-5678-9012-3456",
  "Inhaber":
  {
    "Name": "Mustermann",
    "Vorname": "Max",
    "maennlich": true,
    "Hobbys": ["Reiten", "Golfen", "Lesen"],
    "Alter": 42,
    "Kinder": [],
    "Partner": null
  }
}
\end{verbatim}\normalsize
\end{frame}


%-------------------------------
\begin{frame}[fragile,allowframebreaks]{\LaTeX}
\bg{w3schools}{white}{.tex} Textsatz mit \TeX{} durch Lamport-Makros (= `Shortcuts) \sep
\LaTeX{} liest Makros (mit sprechenden Namen) ein -- in der `Produktion' entsteht die `Illusion' rein deskriptiver Auszeichnung, im Hintergrund sind dies allerdings nur Stellvertreter für die komplexe prozedurale Sprache \TeX . \sep 
prozedural vs. präsentational: \textbf{WYIWYG} (\emph{what you see is what you get}, z.B. Word) vs. \textbf{WYSIWYM} (\emph{what you see is what you mean}) \LaTeX{}-Editoren) \\


\footnotesize
Befehle unten ohne spaces vor der Klammer \\
\mycommand{\textit{Italic}}{\textit{Italic} (darst.)}
\mycommand{\emph{Italic}}{\textit{Italic} (beschr. / semant.)}
\mycommand{\textbf{Bold}}{\textbf{Bold}}
\mycommand{\section{Titel}}{Überschrift 1}
\mycommand{\subsection{Titel}}{usw.}
\mycommand{\href{http://a.com}{Link}}{\href{https://www.latex-project.org/latex3/}{`hidden' Link}}
\mycommand{\includegraphics{bla.png}}{Bild}

\end{frame}

\begin{frame}[standout]
  Übung: Testdokument erstellen auf \alert{\href{https://www.overleaf.com/}{\faCode Overleaf}}
\end{frame}



%-----------------------------------------------------
\begin{frame}[fragile,allowframebreaks]{\href{https://commonmark.org/help}{Markdown}}
\includegraphics[width=0.25\textwidth]{md.png}
\href{https://commonmark.org/help}{Markdown in 60s} \sep simple text formatting ~
\bg{alert}{white}{.md // darstellend}\\

\footnotesize
\mycommand{*Italic*}{\textit{Italic}}
\mycommand{**Bold**}{\textbf{Bold}}
\mycommand{# Heading 1}{Überschrift 1}
\mycommand{## Heading 2}{usw.}
\mycommand{[Link]{http://a.com}}{`versteckter' Link}
\mycommand{![Image][http://url/a.png}{Bild}
\mycommand{> Blockquote}{Zitat}
\mycommand{- List}{Liste (schachtelbar, 2 spaces)}
\mycommand{* List}{Alternative}
\mycommand{1. Aufzaehlung}{Aufzählung}
\mycommand{---}{Trennlinie}
\mycommand{`Inline code`}{Code (`backticks`)}
\mycommand{```code block```}{Codeblock}
\normalsize
\end{frame}

\begin{frame}[standout]
  Übung: \alert{\href{https://commonmark.org/help}{Markdown in 60s}} oder 10min Übung
\end{frame}




\begin{frame}[fragile,allowframebreaks]{\href{https://de.wikipedia.org/wiki/PostScript}{PostScript}}
Seitenbeschreibungssprache \sep 1980er (Adobe Systems) \sep Vektorgraphikformat für Drucker \sep aber auch: Turing-vollständige, stackorientierte Programmiersprache \sep Standard der Druckindustrie \sep heute von PDF (\emph{Portable Dokument Format}) größtenteils abgelöst (ebenfalls Adobe, Weiterentwicklung von PS) \sep kann per PostScript-Druckertreiber aus allen möglichen Dokumenten erstellt werden \sep in Unix per `Ghostscript' verarbeitet \sep Beschreibung von Dokumenten als skalierbare Vektorgraphiken: verlustfreie Vergrößerung 
~\bg{alert}{white}{.ps // präsentational}\\[0.2em]
\footnotesize
Beispielprogramm schreibt `Hallo Welt!' an Postition 50,50. Per Default Koordinatensystem beginnend links unten.
\begin{verbatim}
%!
/Courier findfont    % Schriftart
20 scalefont         % Schriftgröße 20 
setfont              % festlegen
50 50 moveto         % (50, 50)=Schreibposition
(Hallo Welt!) show   % Text ausgeben
showpage             % Seite ausgeben
\end{verbatim}
\normalsize

\end{frame}


\section{Annotation mit XML}
%-----------------------------------------------------
\begin{frame}[allowframebreaks]{XML: eXtensible Markup Language}
\href{https://www.w3schools.com/xml/default.asp}{W3Schools-Tutorial} \\
\bg{alert}{white}{Paradigma der Trennung von Form \& Inhalt} ~
\bg{alert}{white}{XML: Metasprache}\\
\bgupper{w3schools}{black}{.xml}\\
\bg{w3schools}{white}{RSS}~\bg{w3schools}{white}{MathML}~
\bg{w3schools}{white}{GraphML}~\bg{w3schools}{white}{XHTML}~\\
\bg{w3schools}{white}{XAML}~\bg{w3schools}{white}{SOAP}~\bg{w3schools}{white}{RDF}~
\bg{w3schools}{white}{KML}~\bg{w3schools}{white}{OSM}~\\
\bg{w3schools}{white}{Scalable Vector Graphics (SVG)}~(z.B.)

\begin{quote}
    Die \textbf{Extensible Markup Language} (dt. \emph{Erweiterbare Auszeichnungssprache}), abgekürzt XML, ist eine Auszeichnungssprache zur Darstellung hierarchisch strukturierter Daten im Format einer Textdatei, die sowohl von Menschen als auch von Maschinen lesbar ist.

XML wird auch für den plattform- und implementationsunabhängigen Austausch von Daten zwischen Computersystemen eingesetzt, insbesondere über das Internet, und wurde vom World Wide Web Consortium (W3C) am 10. Februar 1998 veröffentlicht. \punkti \textbf{XML ist eine Metasprache,} auf deren Basis durch strukturelle und inhaltliche Einschränkungen anwendungsspezifische Sprachen definiert werden. Diese Einschränkungen werden entweder durch eine \textbf{Document Type Description (DTD)} oder durch ein \textbf{XML Schema} ausgedrückt.  (\href{https://de.wikipedia.org/wiki/Extensible_Markup_Language}{Wiki})
\end{quote}
\end{frame}

%-----------------
\begin{frame}[allowframebreaks]{XML-Familie}
\footnotesize
\bg{w3schools}{white}{XML}~strukturierte Datenbeschreibung \\
\bg{w3schools}{white}{XPath}~Navigation in XMLs \\
\bg{w3schools}{white}{XML Schema}~striktes Datenmodell \\
\bg{w3schools}{white}{XSL}~eXtensible Style Language \\
\bg{w3schools}{white}{XSLT}~Transformation von XML-Dokumenten  \\
\bg{w3schools}{white}{XSL-FO}~ formatierte Ausgabe (z.B. für Druck) \\
\bg{w3schools}{white}{XQuery}~XML-Datenbank-Abfragesprache \\
\bg{w3schools}{white}{und mehr}~
\end{frame}


\begin{frame}[allowframebreaks]{XML-Vokabularien}
\footnotesize
\textbf{(X)HTML} Hypertext Markup Language \sep \textbf{EAD} Encoded Archival Description \sep \textbf{TEI} Text Encoding Initiative \sep \textbf{CEI} Charters Encoding Initiative \sep \textbf{MEI} Music Encoding Initiative \sep \textbf{SVG} Scalable Vector Graphics \sep \textbf{MathML} \sep \textbf{CML} Chemical Markup Language, \dots
\end{frame}

%-----------------------------------------------------
\begin{frame}[allowframebreaks]{XML: eXtensible Markup Language}
\footnotesize
deskriptiv $\neq$ prozedural \sep erweiterbar (nicht wie bei HTML): kein fixiertes Tag-Set \sep 
UTF-8 kodiert

~\bg{alert}{white}{Bedeutung der Daten > Darstellung}\\

expliziert Implizites oder eine Interpretation \sep Hinzufügen von (Meta-)Information \sep maschinelle Weiterverarbeitung 
\sep Standard für Beschreibung und Austausch von Daten ~\bg{alert}{white}{universelle Metasprache}\\


\bg{w3schools}{white}{Nachteile}~  `geschwätzig', Overlapping-Problem, langsam verglichen mit JSON und relationalen Datenbanken -- hat aber auch Features, die die nicht haben (\& für GeWi sehr wichtig sind)

\bg{w3schools}{white}{menschen- und maschinenlesbar}~ 

\bg{w3schools}{white}{unabhängig von Ausgabeformat}~ (Papier, Bildschirm)

\end{frame}

%----------------------------------


\begin{frame}[fragile,allowframebreaks]{XML-Regeln}
\footnotesize
\bg{w3schools}{white}{Prolog}~ \\
\mycommand{<xml version="1.0" encoding="utf-8">}{XML-Deklaration}
\mycommand{<?xsl-stylesheet type="text/xsl" href="mein.xsl"?>}{Verarbeitungsanweisungen (optional)}

evtl. Einbindung des Dokumentmodells (optional) \\
DTD, XML Schema, RelaxNG, Schematron \\
\smallskip

\bg{w3schools}{white}{Entitäten}~ `geschützte Zeichen', da sie eine Bedeutung in der Metasprache und der Objektsprache haben; selbstdefiniert \& vordefinierte, z.B.: \\
\mycommand{&lt;}{<}
\mycommand{&gt;}{>}
\mycommand{&amp;}{\&}
\end{frame}

\begin{frame}[fragile,allowframebreaks]{XML-Regeln}
\footnotesize
Prüfung auf \textbf{Wohlgeformtheit} (nach Regeln des XML-Standards) und \textbf{Validität / Gültigkeit} (wohlgeformt \& Schema-Grammatik entsprechend): kann nur geparst werden, wenn wohlgeformt. Schema kann im Notfall ignoriert werden, auch wenn der Validator sich beschwert.

Es gibt Regeln für Elementnamen. Sollen hier nicht im Detail erklärt werden -- Im Zweifelsfall wird man bei der Validierung erfahren, ob man sich einen Fehltritt geleistet hat.

\includegraphics[width=0.1\textwidth]{doppelkeks.jpg}~\bg{alert}{white}{Doppelkeks} ~ ~\bg{alert}{white}{russische Puppe}~\includegraphics[width=0.15\textwidth]{matroschka.jpg}
\vspace{1em}

\mycommand{<key>value</key>}{XML als key-value-Notation} 

Schachtelung \sep genau 1 Wurzelelement (eine äußerste Puppe)


\bg{w3schools}{white}{Regeln}~ 
Start- und End-Tag \sep Tagnamen case-sensitive \sep leere Elemente erlaubt (\& abkürzbar) \sep ordnungsgemäße hierarchische Schachtelung unter Wurzelknoten

\end{frame}


\begin{mybiggerxml}{Minimalbeispiel}
<?xml version="1.0" ?>
<root>
<element attribute="value">
content
</element>
<!-- comment -->
</root>
\end{mybiggerxml}

%-----------------------------------
\begin{frame}[fragile,allowframebreaks]{Praktische Tipps}
\footnotesize
\bg{w3schools}{white}{Von .docx zu TEI}~ 
Worddokumente (.docx) sind XML-basiert und können, z.B. mithilfe von Pandoc oder OxGarage in einen gewünschten XML-Output umgewandelt werden.
 z.B.~ \href{https://pandoc.org/}{Pandoc} oder \href{http://oxgarage.tei-c.org/#}{OxGarage} (TEI-Fokus)
Mithilfe von XSLT-Transformationen kann man dieses ``Roh-XML'' dann noch anpassen. \\
\smallskip

\bg{w3schools}{white}{Formatvorlagen}~ 
Kursiv, Superskript, fett, etc. geht automatisch, der Rest muss extra angelegt werden. (Formatvorlage = .dotx-Datei aus einem Beispiel-Doc, wo alle Elemente vorkommen, erstellen. Um ein neues Dokument daraus zu erstellen: .dotx klicken). \\
\mycommand{<p rend="lat. Zitat">dixit [...]</p>}{Absatz-Formatvorlagen}
\mycommand{<hi rend="italic">kursiver Text</hi>}{Zeichen-Formatvorlagen}
\normalsize
\end{frame}



\begin{myxml}{`Roh-XML' aus Word mit Formatvorlagen}
<?xml version="1.0"?>
<xml>
    <p rend="la::Dichtung">
        <hi rend="la::Struktur-Heading">Schluss;;Epimythion::11</hi></p>
    <p rend="la::Dichtung">
        <hi rend="metrAnalyse">#senar::u-|u-|u//-|u/-|u-|ux</hi>
        sic totam praedam sola improbitas
        <note place="foot" xml:id="ftn15" n="15">
            <p rend="footnote text">
                Vokabel::improbitas,-atis f.: Schlechtigkeit (siehe auch. improbus)
            </p>
        </note>
        abstulit. 
    </p>
    <p rend="de">
        <hi rend="de::Title">Die Kuh, die Ziege, das Schaf und der Löwe</hi>
    </p>
</xml>
\end{myxml}


\section{Text Encoding Initiative}
%-----------------------------------------------------
\begin{frame}[fragile,allowframebreaks]{TEI Primer}
\footnotesize
\bg{alert}{white}{Text Encoding Initiative}\\
\bgupper{w3schools}{black}{.xml}\\
XML-Standard \sep De-Facto-Standard in den \emph{Digital Humanities} \sep Kodierung von Druckwerken (Editionswissenschaft) \sep bishin zu Linguistik
\begin{quote}
    Die Text Encoding Initiative (TEI) ist eine 1987 gegründete Organisation (seit 2000 als TEI-Konsortium organisiert) und ein gleichnamiges Dokumentenformat zur Kodierung und zum Austausch von Texten, das diese entwickelt hat und weiterentwickelt.  (\href{https://de.wikipedia.org/wiki/Text_Encoding_Initiative}{Wiki})
\end{quote}

\begin{myxml}{TEI-Grundstruktur}
<TEI> <!-- Wurzelelement -->
    <teiHeader> ... </teiHeader> <!-- Autor, Titel, Datierung, Quelle(n), Editionsrichtlinien,
Versionierung, etc. -->
    <text> ... </text>
</TEI>
\end{myxml}
\smallskip

\bg{w3schools}{white}{Ressourcen}\\
\href{http://www.tei-c.org/Support/Learn/}{Learn TEI} \sep \href{http://www.tei-c.org/support/learn/teach-yourself-tei/}{Teach Yourself} \sep P5 = 5. Proposal \sep MEI für Musik \sep CEI für Charters (Urkunden) \sep \href{http://www.tei-c.org/}{http://www.tei-c.org/} \sep Standard zur Textkodierung \sep Publikationswerkzeuge: Versioning Machine, TEI Boilerplate, TEICHI, TEI
Stylesheets 
\end{frame}

%-------------------------------------
\begin{frame}[fragile,allowframebreaks]{TEI Header}
\footnotesize
\bg{alert}{white}{fileDesc}~ = Instrumentarium zur umfassenden bibliographischen  Beschreibung der Inhalte eines TEI-Dokumentes  \\
\bg{alert}{white}{encodingDesc}~ = beschreibt den Zusammenhang des elektronischen Textes mit
dem Quelltext, z.B. diverse Regel bei der Transkription, sowie Erläuterungen zum
Annotationsprozess \\

\begin{myxml}{TEI-Header}
<TEI> <!-- Wurzelelement -->
    <teiHeader>
        <fileDesc> ... </fileDesc> <!-- obligatorisch --> 
        <encodingDesc> <!-- optional -->
        <profileDesc> <!-- optional -->
        <revisionDesc> <!-- optional -->
    </teiHeader> 
    <text> ... </text>
</TEI>
\end{myxml} 
\bg{alert}{white}{profileDesc}~ = Beschreibung aller nicht bibliographischen Aspekte des
Textes, z.B. Informationen über die Entstehung des Textes, sowie über verwendete Sprachen
u.Ä. \\
\bg{alert}{white}{revisionDesc}~ = Beschreibung aller Änderungen und Überarbeitungsschritte am
Transskript des Quelltextes.
\end{frame}
%----------------------------------
\begin{frame}[allowframebreaks]{TEI Basics}
\footnotesize
TEI = \bg{alert}{white}{modular} $\to$ im Schema können Untermengen festgelegt werden (`Ich benutze Core und \dots ') \sep \href{http://www.tei-c.org/Roma/}{ROMA Schema} \sep sonst TEI P5 All
\smallskip

\bg{w3schools}{white}{Gemeinsamkeiten von Textdokumenten: } \\
\bg{alert}{white}{\href{http://www.tei-c.org/release/doc/tei-p5-doc/en/html/CO.html}{TEI Core} }

\begin{itemize}\footnotesize
    \item Identifikatoren (Seitenangaben, Signatur, Inventarnummer, Regalnummer, etc.)
    \item Abschnitte und Unterabschnitte (Gliederung)
    \item Abbildungen, Skizzen, grafische Elemente
    \item Schreibmodus (Prosa, Drama, Vers etc.)
    \item Strukturelle Einheiten (Absatz, Listen, Strophen, Zeilen, Reden, etc.)
    \item Textunterschiede oftmals durch unterschiedliche Formatierung gekennzeichnet (Titel, Überschriften, Zitate, Betonungen, etc.)
    \item Sachinformationen (Personen, Orte, etc.)
    \item Texteingriffe Korrekturen, Streichungen, Revisionen
\end{itemize}
\vspace{1em}

\bg{w3schools}{white}{Warum TEI verwenden oder nicht?} \\
\textbf{Nachteil:} man muss Regeln befolgen -- muss man aber sonst sinnvollerweise eh auch. Wollte man es wirklich gut machen, wäre es viel Arbeit! Andererseits: Flexibilität. TEI hat nicht immer etwas für Spezialzwecke. \\
\textbf{Vorteil:} die Regeln hat sich schon jemand überlegt, man muss sie nur nachschauen, verstehen und anwenden. Allgemeiner Standard in den Geisteswissenschaften, macht Daten unterschiedlicher Projekte interoperabel.
\end{frame}


%-----------------------------------
\begin{frame}[fragile,allowframebreaks]{TEI verwenden}
\footnotesize
\href{http://www.tei-c.org/release/doc/tei-p5-doc/en/html/SG.html}{Gentle Intro to XML}

\bg{alert}{white}{TEI Core}~ \textbf{div} (Abschnitt) \sep \textbf{p} (Paragraph) \sep \textbf{head} (Überschrift) \sep \textbf{lb} (linebreak) \sep \textbf{pb} (page break / beginning) \sep \textbf{hi} (highlight) \sep \textbf{l} (line) \sep \textbf{lg} (line group) \sep \textbf{list} \sep \textbf{item} \sep \textbf{listBibl} \sep \textbf{bibl} (bibliograph. Angabe)

\bg{alert}{white}{Attribute}~ \textbf{@n} (label) \sep \textbf{@type} (Typisierung) \sep \textbf{xml:id} (eindeutige Kennung) \sep \textbf{xml:lang} (Sprache) \sep \textbf{@rend} (Darstellung) \sep \text{@ana} (Interpretation)

\begin{myxml}{Typische Anwendungen}
<foreign xml:lang="en">word</foreign>
<term type="homonym"/>
<date when="2009-04-27"/>
<time when="12:00:00"/>
<name type="person"/>
<persName n="Caesar" xml:id="#44BC">Caesaris</persName> <!-- oder -->
<persName key="ID.01.208"/>
<person/>
<emph/> <hi rend="italic">kursiver Text</hi>
<seg/> <abbr type="acronym"/>
<placeName xml:id="#Whitby">Abbey</placeName>
\end{myxml}

\bg{alert}{white}{Namensräume}~ Identifikation über URI \sep \bg{w3schools}{white}{<präfix:name>}~ \sep z.B. <tei:p> (`Ich meine das <p> nach dem TEI-Standard'). \\
\bg{alert}{white}{Deklaration}~ <element xmlns=“URI“> \dots \\
<präfix:element xmlns:präfix=“URI“> \dots \\
Z.B. <tei:p xmlns:tei=“http://www.tei-c.org/ns/1.0“>...


\end{frame}



\begin{frame}[fragile,allowframebreaks]{Spezialfälle: Rede \& Briefe}
\footnotesize
Kodierung von Sprechakten  (\href{http://www.tei-c.org/release/doc/tei-p5-doc/en/html/examples-sp.html}{TEI-Referenz}), falls Speaker davor steht, sonst \href{http://www.tei-c.org/release/doc/tei-p5-doc/en/html/examples-said.html}{TEIs `said'}:
\begin{verbatim}
<sp who="#person">
    <speaker>1.</speaker> <p>Bla, bla, bla.</p>
</sp>

<said who="#Adolphe">- Alors, Albert, quoi de neuf?</said>
\end{verbatim}

\begin{myxml}{Briefe in TEI \href{http://www.tei-c.org/release/doc/tei-p5-doc/en/html/DS.html#DSOC}{(Referenz)}}
<div type="letter" n="14">
    <head>Letter XIV: Miss Clarissa Harlowe to Miss Howe</head>
        <opener>
            <dateline>Thursday evening, March 2.</dateline>
            <salute>Hallo,</salute>
        </opener>
    <p>On Hannah's depositing my long letter ...</p>
    <closer>
        <salute>Yours more than my own,</salute>
        <signed>Clarissa Harlowe</signed>
    </closer>
</div>
\end{myxml}
\end{frame}

\begin{frame}[fragile,allowframebreaks]{Annotation fortgeschritten}
\footnotesize
\bg{w3schools}{white}{\emph{Named Entities} \& Indirekte Referenz}\\
\href{http://www.tei-c.org/release/doc/tei-p5-doc/en/html/ND.html}{TEI 13: Names, Dates, People, Places} \sep 
\textbf{persName} für Namensnennung, \textbf{<rs>} für \emph{referring string} bei indirekter Nennung (`er', `der Herr', etc.). Hier dann mit \textbf{@key} oder \textbf{@ref} spezifizieren, wer gemeint ist. (\href{http://www.tei-c.org/release/doc/tei-p5-doc/en/html/CO.html#CONARS}{Referenz}). \textbf{forename} \sep \textbf{surname} \seo \textbf{roleName} (z.B. `König') \sep  \textbf{genName} (`der Ältere') \sep \textbf{addName} \sep \textbf{nameLink} (`von').

\begin{verbatim}
<name role="writer" type="person"
ref="http://d-nb.info/gnd/118540238">
Goethe</name>
<person>
  <addName type="Former">Murray</addName>
  <forename>Wilhelmina</forename>
  <addName type="nickname">Mina</addName>
</person
\end{verbatim}

Eine \textbf{person} (selbst) ist nicht identisch mit dem \textbf{persName} (Name)! Dafür zur Verfügung stehende Elemente: \textbf{persName} zu \textbf{person}, \textbf{orgName} für \textbf{org} (Organisation), \textbf{placeName} für \textbf{place.} \textbf{geogName} (= geographical) = Lanschaftsmarker, Berge, etc.
Problem: \textbf{Unterscheidliche Namensformen}, daher:
\bg{w3schools}{white}{Normalisierung:} Zeigen, dass \emph{dieselbe} Person gemeint ist. Wir können also eine sog. `normalisierte' Form angeben, z.B. in Form einer Referenz auf die Personenliste im TEI-Header, auf deren \textbf{@xml:id} wir per \textbf{@ref} verweisen. Oder Normdaten der GND oder in Attribut (z.B. \textbf{@n} (\emph{label}) oder \textbf{@ana} (Interpretation)). 

\bg{w3schools}{white}{Redundanz ist eine Fehlerquelle} $\to$ Referenz auf eindeutigen Ort, wo die Info genau 1x vorhanden ist $\to$ Bei Fehlern nur 1x ausgebessern, nicht 200 Okkurenzen. Zur eindeutigen Referenzierung xml:ids, z.B. \verb|<person xml:id="Mina">Mina</person>|, darauf refernzieren dann mithilfe des \texttt{@ref}-Attributs und einem Hashtag: \verb|<persName ref="#Mina">Mina</persName>|.

\end{frame}


%---------------------------------------------
\section{Annotationspraxis}
\begin{frame}[allowframebreaks]{Hands-On gezielte Annotation}
\bg{alert}{white}{= Annotation mit dem Ziel der Erschließung}\\ 
\bg{w3schools}{white}{Vorgehensweise}
\smallskip 

 \footnotesize
\begin{enumerate}
    \item Forschungsfrage formulieren
    \item Überlegen, welche Daten zur Beantwortung benötigt werden
    \item Begründen, warum und wie diese Daten zur Beantwortung beitragen
    \item gezielt zu erschließende Quelle / Objekt (kann auch ein Gegenstand sein!) annotieren
\end{enumerate}

Dazu ist (wahrscheinlich) \bg{alert}{white}{Vorbereitung} notwendig:\\[0.5em]
\end{frame}

\begin{frame}[allowframebreaks]{Beschaffung der Daten}

\begin{enumerate}
    \item Woher kann ich die Daten bekommen? 
    \item Online in welcher Form (xml, html, txt, etc.) verfügbar?
    \item Falls ja: darf ich die Daten (für meine Zwecke) überhaupt verwenden? (Urheberrecht, Creative Commons, etc.)
    \item selbst digitalisieren / transkribieren
    \item gezielt zu erschließende Quelle / Objekt (kann auch ein Gegenstand sein!) annotieren
\end{enumerate}

\end{frame}

\begin{frame}[allowframebreaks]{Sog. Pre-Processing}

 Kann je nach Daten und Bedürfnisse ganz unterschiedlich aussehen. \\
 
\begin{enumerate}
    \item `Junk' / `Noise', der entfernt werden soll? (Zeichen aus dem HTML, kaputtes Encoding, irrelevante Daten, etc.)
    \item Basis-Kodierung könnte man unter Pre-Processing fassen (Überschriften, Absätze)
    \item Können Tätigkeiten automatisiert werden? Gibt es die Ressource schon online im richtigen Format?
    \item Aber Achtung, übermotiviertes Automatisieren kann in Mehrarbeit enden!
    \item Überprüfen, ob die `Vereinfachung' tatsächlich das gewünschte Ergebnis erzielt!
    \item Falls Online-Markup nicht passt: Evtl. ist es mit einem plain text Dokument im Endeffekt schneller. 
    \item Advanced: Automatisierung / Pre-Processing mit Python, XSLT, etc.
\end{enumerate}

\end{frame}
%----------------------------------

\begin{frame}[allowframebreaks]{Annotation leichter machen}
\footnotesize
\begin{enumerate}
    \item beim Annotieren nicht `lesen', für das Grobe nur überfliegen
    \item erst das Grobe (ohne zu lesen), dann eine 2. Runde für Details (\emph{Named Entities}, etc.)
    \item Keyboard-Shortcuts nutzen: per Tastatur navigieren lernen statt `herumklicken'
    \item Suchen \& Ersetzen für häufige Orte / Namen
    \item Patterns automatisch finden (RegEx, Suchen/Ersetzen)
    \item vielleicht gibt es die Daten schon online??
\end{enumerate}
\smallskip

Wir könnten alles in Word mit Formatvorlagen annotieren und per OxGarage umwandeln. Allerdings muss dort (bis auf Absätze) alles händisch markiert werden. Suchen \& Ersetzen wäre da wohl effizienter.

\begin{enumerate}
    \item \textbf{Basiskodierung I:} Daten aus Wikisource rauskopieren in Word. Alles markieren (Strg+A) und auf Formatvorlage `Standard' klicken. Testen, ob es wirklich übernommen wurde. Im Word dann mit ``Überschrift''-Formatvorlage die Kapitelüberschriften markieren. Würde auch mit Markdown schnell gehen. 
    (z.B. jedes `Chapter' durch `\# Chapter' ersetzen $\to$ in OxGarage zu TEI machen lassen, dann sind die \textbf{heads} schon drin (Achtung, dazu muss man das \texttt{.txt} in eine \texttt{.md}-Endung umbenennen, sonst wird es nicht akzeptiert.
    \item \textbf{Transformation zu XML} / TEI mithilfe von \href{http://oxgarage.tei-c.org/#}{OxGarage}.
    \item \textbf{Tipp:} Wenn ein Schema (z.B. TEI) mit dem Dokument verbunden ist, einfach öffnende Spitzklammer an eine Stelle schreiben und Oxygen bietet an, was dort stehen könnte (z.B. beim Header nützlich). Geht auch für Attribute, aber geht natürlich nur, wenn man mit dem Internet verbunden ist. Wenn man z.B. <p schreibt, dann zeigt es alle hier möglichen Elementnamen beginnend mit p, falls man schon ungefähr weiß, dass man z.B. eine Person sucht, die so benannt sein könnte.
    \item \textbf{Suchen \& Ersetzen nutzen,} um häufige Orte und Personennamen `mit einem Schlag' zu annotieren. Sicherstellen, dass man nicht Dinge markiert, die keine eigentlichen Treffer waren (\emph{false positives}): Zunächst ein paar Ergebnisse durchklicken, bevor man `Alles ersetzen' macht. Vor Namen sicherheitshalber ein Leerzeichen setzen, damit nicht gleiche Wortteile mitgenommen werden (z.B. bei Mina durchaus möglich). Wenn Leute mit Vor- und Nachnamen vorkommen, immer zuerst die längstmögliche Version suchen, dann z.B. auch den Vornamen ohne Nachnamen. Allgemein, auch wenn man Suchen \& Ersetzen zum Bereinigen unerwünschter Attribute verwendet, immer zuerst das längstmögliche, dann die kürzeren Varianten. Beim `Bereinigen' schauen, ob man die unerwünschten Attribute nicht eh vielleicht brauchen kann (\emph{<hi style="bold">} könnte z.B. \emph{<hi rend="Name">} werden). (Für RegEx-Experten: XSLT / oXygen unterstützt nicht alle RegEx, die ihr evtl. aus anderen Programmiersprachen kennt -- hier sorgfältig googeln: Falls es den Eindruck macht, dass der gewünschte Ausdruck nicht explizit irgendwo als XSLT-tauglich ausgewiesen wird, davon ausgehen, dass es ihn nicht gibt. Mangelnde Bereitschaft zur Akzeptanz dieses Umstands kann in stundenlanger erfolgloser Fehlersuche enden \faSmileO )
    \item Text durchgehen auf \emph{Named Entities} (Namen, Orte), die nicht automatisch gefunden wurden. Per Hand nachziehen.
\end{enumerate}

\end{frame}


\begin{frame}[allowframebreaks]{Vorgehensweise zur Basiskodierung}
\footnotesize
\begin{enumerate}
    \item Annotation mit Word-Formatvorlagen und OxGarage-Konversion
    \item XML selbst schreiben mit Tastatur-Shortcuts
    \item im Internet vor-kodierten Text suchen / aus XHTML transformieren, etc. (funktionierte hier nicht so gut)
\end{enumerate}
\smallskip

\bg{alert}{white}{Dateien}~sinnvolle, aussagekräftige Dateinamen. Keine Leerzeichen in Dateinamen, keine Umlaute. Möglichst plattformunabhängige Dateiformate wählen.


\end{frame}

\begin{frame}[allowframebreaks]{Basis-Kodierung}

\begin{enumerate}
    \item \textbf{Überschriften} (\emph{head}) und \textbf{Paragraphen} (\emph{p})
    \item händisch erstellter, \textbf{minimaler TEI-Header}
    \item \textbf{Personennamen, Orte} (mindestens direkte Verwendung, Verweise wie `er' optional)
    \item \textbf{direkte Reden mit eindeutigem Attribut} (\emph{listPerson} im \emph{teiHeader}), das angibt, wer hier spricht
    \item \textbf{Briefe} (ganzer Brief, als Unter-Verschachtelung: Absender, Empfänger, Ort, Datum nach TEI)
\end{enumerate}

\bg{alert}{white}{TEI-Header}
\begin{enumerate}
    \item \textbf{fileDesc} ist verpflichtend, hier einfach die Autorschaft (ihr selbst und z.B. Bram Stoker) klarstellen und ggf. erklären
    \item dann in der \textbf{profileDesc}: \textbf{listPerson} $\to$ \textbf{person} (mit Attribut xml:id="Mina" oder so)  $\to$ persName, evtl mit Untergliederung darunter nochmals in Vor-, Nach- und Spitzname. Für alle Hauptfiguren. TEI-Konventionen: \href{http://www.tei-c.org/release/doc/tei-p5-doc/de/html/examples-addName.html}{addName} type="nick" für Spitznamen, \textbf{roleName} für Graf, Doktor etc.
    \item im Text Personen dann mit \texttt{<persName ref="#Mina">} referenzieren (selbes entsprechend bei Orten)
    \item Liste der wichtigsten Orte: \textbf{profileDesc} $\to$ \textbf{settingDesc} $\to$ \textbf{listPlace}: \texttt{<place xml:id="London"  ana="http://www.geonames.org/2643741/"> <placeName> London} etc.
    \item im Text: \texttt{<placeName ref="#London">London</placeName>}
\end{enumerate}

\end{frame}

%-----------------------------------------------------
%


\begin{frame}[standout]
  \alert{Und Action!} \normalsize
 Erstellt einen TEI-Header für Euer Dokument.
\end{frame}


%---------------------------------------------
\begin{multibox}{2} % Anzahl der Boxen in einer Reihe angeben
\begin{subbox}{subbox}{Tastatur-Shortcuts}\footnotesize

\mycommand{STRG+C}{copy}
\mycommand{STRG+Z}{undo}
\mycommand{STRG+Y}{redo}
\mycommand{STRG+A}{alles}
\mycommand{STRG+S}{save}

\mycommand{Alt+Tab}{zw. Anwendungen springen}
\mycommand{Strg+Tab}{zw. Tabs springen}
\mycommand{F5}{Browser-Refresh}

\end{subbox}
\begin{subbox}{customcolor}{Tastatur-Shortcuts}\footnotesize

\mycommand{STRG+Pfeiltasten}{in Sprüngen navigieren}
\mycommand{STRG+SHIFT+Pfeiltasten}{in Sprüngen markieren)}
\mycommand{STRG+SHIFT+e}{Element (<oXygen/>)}

\end{subbox}
\end{multibox}

\begin{frame}[fragile,allowframebreaks]{RegEx}
\footnotesize

\href{http://regexr.com}{RegExr.com}: vorher testen (potentielle Probleme z.B. = entweder alles oder nichts finden).
Auch ein RegEx ist ein Modell $\to$ Abstraktion. z.B. \textbf{Jahreszahl finden} \verb|\d{4}|
\smallskip

\textbf{Ein Wort finden, das auch als Wortteil woanders vorkommen kann}, z.B. `bat': Suche von \verb|bat(\W)|. D.h. die Zeichenkette beginnt mit `bat' und es folgt ein Nicht-Wortzeiten (also Space, Punkt, Komma, Hyphen, etc.). \verb|\w| gibt ein `Wortzeichen', also Buchstaben, der Ausdruck in Groß jeweils das Gegenteil -- also alles, was durch \verb|\w| \emph{nicht} erfasst wird.
Das ganze kann man natürlich auch an den Wortanfang noch setzen, um auf Nummer sicher zu gehen: \verb|(\W)bat(\W)|
\smallskip

Reguläre Ausdrücke (Regular Expressions, Regex) = Muster, die Zeichenketten beschreiben. 
$\to$ Finden und Bearbeiten von Zeichenfolgen, die auf sie zutreffen.

\begin{multibox}{2} % Anzahl der Boxen in einer Reihe angeben
\begin{subbox}{subbox}{nachgestellt}\footnotesize

\mycommand{nichts}{genau 1x}
\mycommand{?}{1x oder kein Mal}
\mycommand{*}{beliebig oft}
\mycommand{+}{mind. 1x}
\mycommand{{n}}{n-mal}

\end{subbox}
\begin{subbox}{customcolor}{Auswahl}\footnotesize

\mycommand{.}{beliebiges Zeichen}

\mycommand{[abc]}{Zeichenauswahl}
\mycommand{[a-z][0-9]}{Auswahl}
\mycommand{\n}{Zeilenumbruch}

\end{subbox}
\end{multibox}

%---------------------------------------------
\begin{multibox}{2} % Anzahl der Boxen in einer Reihe angeben
\begin{subbox}{subbox}{Gruppieren}\footnotesize

\mycommand{()}{zur Gruppierung und Anwahl (\$1)}
\mycommand{(.*?)}{beliebig (non-greedy)}


\mycommand{^}{negiert Folgendes}
\mycommand{^}{Start eines Strings}
\mycommand{\$}{String-Ende}
\mycommand{\\}{Escape-Sequenz}
\mycommand{|}{Oder-Pipe}
\end{subbox}
\begin{subbox}{subbox}{RegEx}\footnotesize
\mycommand{\s}{\emph{space}: Leerzeichen, Tab, Zeilenumbruch}
\mycommand{\d}{Zahl (\emph{digit}}
\mycommand{\w}{Wort (=Buchst.,Zahl, Unterstrich}
\mycommand{\D \W \S}{jeweils das Gegenteil wie in klein}

\end{subbox}
\end{multibox}

\end{frame}






\begin{frame}[fragile,allowframebreaks]{Reguläre Ausdrücke (RegEX)}
\footnotesize
\bg{w3schools}{white}{Vorbereiten}\\
Als erstes: Ein Leerzeichen markieren und durch ein Leerzeichen ersetzen. Dies verhindert Encoding-Probleme (z.B. Windows hat `geschützte Leerzeichen', die im Hintergrund für den Computer anders aussehen. Diese können unerklärliche Fehler verursachen, daher lieber gleich am Anfang präventiv erledigen und die Fehlerquelle damit ausschließen.
\medskip


\bg{w3schools}{white}{Personen und Orte}\\
\textbf{Suchen und Ersetzen in Oxygen:} Im OxygenXML-Suchfeld `Reguläre Ausdrücke' und alles drunter anklicken.

Im oberen Suche-Feld:
\begin{verbatim}
    ( Mina|Mina Murray|Mina Harker)
\end{verbatim}

Im Ersetzen-Feld greifen wir auf die Inhalte aus der Klammer folgendermaßen zu (mit \$, das durchnummeriert wird, denn man kann mehrere Ausdrücke klammern und somit unterschiedlich weriterverarbeiten (z.B. Vor- und Nachname bereits richtig kodieren mit (Mina)(Murray|Harker), der Vorname ist dann in \$1, der Nachname in \$2 wiederzufinden):
\begin{verbatim}
    <persName ref="#Mina">$1</persName>
\end{verbatim}

Selbes für Ort:
\begin{verbatim}
    <placeName ref="#London">$1</placeName>
\end{verbatim}

Hier nicht vergessen, dass im Text immer vor der id der \# sein muss, sonst kommt vom TEI-Schema ein Fehler gemeldet.

\medskip

\bg{w3schools}{white}{Kapitelstruktur}\\
Im oberen Suche-Feld:
\begin{verbatim}
    <p>(Letter|Chapter) (.*?)</p>
\end{verbatim}

Im Ersetzen-Feld greifen wir auf die Inhalte aus der Klammer folgendermaßen zu (mit \$):
\begin{verbatim}
    <head>$1 $2</head>
\end{verbatim}

Hinterher muss man um die heads (also immer beginnend vor head markieren bis kurz vor dem folgenden head, d.h. Ende Kapitel) markieren, Strg+E und ein div-Element erstellen für jedes Kapitel. Sonst nicht TEI-konform (gäbe auch komplexe Abfragen, mit denen dies automatisiert werden könnte).
\medskip

\bg{w3schools}{white}{Unerwünschte Zeilenumbrüche loswerden}\\
Z.B. bei den Projekt-Gutenberg-txt-Dateien ist es uns passiert, dass unerwünschte Zeilenumbrüche im Text waren, die ohne RegEx nicht so einfach zu entfernen waren. Der RegEx, der das Problem löst, schaut folgendermaßen aus:
\begin{verbatim}
    (\w+)(\n)(\w+)
\end{verbatim}
D.h. er sucht ein oder mehrere Wörter, wo genau 1 Zeilenumbrucht (\textbackslash n = newline) dazwischen steht. Die Logik dahinter: richtige Absätze werden durch mehrere Zeilenumbrüche abgebildet (bei einer Leerzeile zw. zwei Abschnitten finden sich mind. 2 newline-Zeichen, also 1x am Ende der letzten Zeile und 1x nach der Leerzeile, dort sind aber keine Wort-Zeichen dazwischen. Wir löschen also tatsächlich nur unerwünschte / überflüssige  Umbrüche. Hier aber auch sichergehen, dass der Text tatsächlich so `funktioniert' -- wenn keine Leerzeilen zwischen Abschnitten sind, so funktioniert das nicht!

\end{frame}
%---------------------------------------------------



\begin{frame}[allowframebreaks]{Literatur}

\AtNextBibliography{\footnotesize}
\printbibliography  
\end{frame}








%--------------------------------------------------

\section{Webentwicklung}


\begin{frame}[fragile,allowframebreaks]{HTML}

\bg{alert}{white}{HTML = Hyper Text Markup Language}\\
\bgupper{w3schools}{black}{.html}\\
Struktur von Webseiten. Ähnliche wie XML, aber fixiertes Set an Tags (viel weniger, die man sich merken muss und oft braucht).

Die wichtigsten HTML-Elemente: 
html, head, body.
div, p, h1-6.
span,
ul, ol, li.
table, tr, td.
%------------------
\smallskip

\begin{myhtml}{HTML-Grundstruktur}
<!DOCTYPE html>
<html>
<head>
<title>Page Title</title>

<style> /* inline CSS, lieber in extra file */
h1 {
  color: blue;
  font-family: verdana;
  font-size: 300%;
}
p  {
  color: red;
  font-family: courier;
  font-size: 160%;
}
p.important {
    color: green;
}
</style>
<script>
alert("Hello! I am an alert box!!");
</script>
</head>
<body>

<!
<h1>This is a Heading</h1>
<p>This is a paragraph. <br />
<a href="https://www.w3schools.com">This is a link</a>
</p>

<img src="img.jpg" width="500" height="600" />
<p style="color:red">I am a paragraph</p>

<p title="I'm a tooltip">
This is a paragraph.
</p>
<p>My mother has <span style="color:blue">blue</span> eyes.</p>
<p class="important">Note that this is an important paragraph. :)</p>


</body>
</html>
\end{myhtml}
\end{frame}

\begin{frame}[fragile,allowframebreaks]{CSS}
\bg{alert}{white}{CSS = Cascading Style Sheets}\\
\bgupper{w3schools}{black}{.css}\\

wird per Link in HTML eingebunden \sep beschreibt die Darstellung von Webseiten \sep
Trennung von Inhalt und Darstellung: HTML hat den Inhalt in strukturierter Form, CSS macht das Layout.

Das \bg{w3schools}{white}{\href{http://getbootstrap.com/}{Bootstrap-Framework}} bietet schon viele fertig anzuwendende Elemente. $\to$ Man muss nicht alles von Hand machen, daher sehr empfehlenswert. \sep Box- und Grid-Modell

\href{https://tutorialzine.com/2015/10/learn-the-bootstrap-grid-in-15-minutes}{15min to Bootstrap}

CSS 3 Cheat Sheet
\href{http://www.smashingmagazine.com/wp-
content/uploads/images/css3-cheat-sheet/css3-cheat-
sheet.pdf}{CSS3-Cheatsheet}
• W3 School
\href{http://www.w3schools.com/css/}{W3schools CSS} \sep \href{http://www.w3.org/TR/CSS21/propidx.html}{Tabelle der Eigenschaften}


\href{http://www.csszengarden.com}{CSS ZenGarden}

\bg{alert}{white}{CSS Syntax}\\

\begin{mycss}{Selektor}
h1 {
font-family: Arial
}
\end{mycss}


\begin{mycss}{Deklarationsblock}
p {
font-family: Arial ;
color: red
}
\end{mycss}

\begin{mycss}{Schrift für \texttt{class}}
.person {
    font-weight:bold;
    font-size:smaller;
    font-style:italic;
}
\end{mycss}
\end{frame}

\begin{frame}[fragile,allowframebreaks]{Bootstrap-Framework}
Von diesem \href{https://getbootstrap.com/docs/4.0/getting-started/introduction/}{Link zu Bootstrap} müssen ganz oben im \texttt{<head>} des HTML einerseits das Bootstrap-Stylesheet, andererseits JQuery-Plugins hineingeladen werden.
(Achtung, \texttt{head} ist hier das Äquivalent zum \texttt{teiHeader}, Überschriften heißen in HTML \texttt{h1-h6}.
Bei \href{https://getbootstrap.com/docs/4.0/examples/}{Examples} gibt es Beispielseiten, die man als Basis verwenden und anpassen kann.
\begin{myhtml}{html head (verkürzt)}
<html>
    <head>
    <meta charset="utf-8">
    <link rel="stylesheet" href="https://...bootstrap.min.css">
    <script src="https://code.jquery.com/jquery-3.2.1.slim.min.js"></script>
    </head>
    <body> [...] </body>
</html>
\end{myhtml}
\href{view-source:https://getbootstrap.com/docs/4.0/examples/starter-template/}{Source-Code des Starter Template} (bei normalem Link: Rechtsklick `Element untersuchen').
\end{frame}

\begin{frame}[fragile,allowframebreaks]{JavaScript}
\bg{alert}{white}{Dynamische Veränderung von Webseiten}\\
\bgupper{w3schools}{black}{.js}\\
JS kann Darstellung in Webseiten verändern, ohne dass man neu laden müsste (daher `dynamisch').

Wenn wir z.B. Dinge per Checkbox wegklicken oder zuschalten wollen, verwenden wir JavaScript.

\begin{myjs}{Hello World}
document.getElementById("demo").innerHTML = "Hello JavaScript";
\end{myjs}

\end{frame}

\begin{frame}[allowframebreaks]{Web-Entwicklung-Intro}
\scriptsize
\bg{alert}{white}{Ressourcen}
\href{https://dash.generalassemb.ly/projects}{Lektionen 1 und Beginn von 2 auf Dash (Anmeldung nötig)}. \sep 
\href{https://www.codecademy.com/learn/learn-html}{Codecademy HTML} (gut, aber eher ausführlich)


Interneting is hard:
\href{https://internetingishard.com/html-and-css/introduction/}{Intro} \sep
\href{https://internetingishard.com/html-and-css/basic-web-pages/}{Kapitel 2: HTML Basics}
\sep 
\href{https://developer.mozilla.org/en-US/docs/Learn/HTML/Introduction_to_HTML/Getting_started}{Mozilla Learn HTML} \& \href{https://developer.mozilla.org/de/docs/Learn/HTML}{auf Deutsch}
\sep
\href{https://www.w3schools.com/html/exercise.asp}{W3Schools interaktives HTML} oder wahlweise textlastiges \href{https://www.w3schools.com/html/default.asp}{HTML-Tutorial}. \sep
\href{https://www.learn-html.org/en/Basic_Elements}{Learn HTML.org} \sep
\href{https://jgthms.com/web-design-in-4-minutes/}{Web Design in 4 minutes (CSS)}
\smallskip


\bg{alert}{white}{Responsive Web Design}
\href{https://www.w3schools.com/html/html_responsive.asp}{W3Schools-Tutorials}

\end{frame}




\begin{frame}[standout]
  \alert{Und Action!} \normalsize
 Wählt aus den Links der letzten Folien ein Tutorial aus und probiert es durch.
\end{frame}


%-----------------------------------------------------

\section{Datenmodellierung}
\begin{frame}[allowframebreaks]{Datenmodellierung}
\bg{w3schools}{white}{Model = Abbildung = Repräsenation $\neq$ Original}, sondern nur ausgewählte Teile des Originals werden durch es wiedergegeben bzw. dargestellt. Die Wahl dieser abzubildenden Aspekte der Realität macht die Nützlichkeit oder Repräsentativität für einen bestimmten Anwendungsfall aus.

\bg{w3schools}{white}{Subjektiv \& abstrahiert} $\to$ Muss auf eine spezielle Frage hin erstellt werden und nützt auch nur für diese. Wahrnehmung und Darstellung sind immer selektiv, so auch die Modelle. Ich kann nie alle Aspekte einer Vase gleichzeitig aufnehmen. Ein Foto zeigt sie nur von einer Seite.

\bg{w3schools}{white}{Metadaten} können zur Anreicherung von Datenmodellen gesammelt werden (Zusatz-Informationen über etwas, administrative \& technische).

\bg{w3schools}{white}{Disambiguierung notwendig}, z.B. Normalisierung, Geonames, GND

\bg{w3schools}{white}{Modell = Interpreatation} Alles ab Basiskodierung, bereits teils bei \emph{Named Enitites}. \textbf{Vorschlag zur Kenntlichmachung dieses Umstands:} \texttt{@ana} verwenden.

Wo endet die `Abbilung von Realität', ab wo ist es nur noch eine Perspektive darauf? 

%Konzeptuelles vs. logisches Datenmodell. Entitäten und Relationen (\emph{Entity Relationship Modell} - konzeptuell

\bg{w3schools}{white}{Was ist für uns überhaupt das Original?} Das Buch, dessen Materialität im Digitalen abstrahiert wird? Die Idee?


\bg{alert}{white}{Modelle = vereinfachte Repräsentationen von Teilen der realen Welt.}

Bsp. Photogrammetrie / Structure from Motion (= aus Fotos 3D-Modelle erstellen): Realität $\to$ mehr Info als in den Bildern. 3D-Modell: es braucht ca. 8 Bildpixel für 1 Punkt im Modell $\to$ weniger Information als die Summe der (mind. 50) Bilder.
\end{frame}

\begin{frame}[allowframebreaks]{Wechseln zw. Modellen}
\bg{w3schools}{white}{Nochmals zu Trennung von Inhalt und Darstellung} Warum nicht gleich in HTML? $\to$ XML zu HTML ist Reduktion, da HTML ein ein viel beschränkteres Modell ist! Daher in den  `abstrakten' Rohdaten den Inhalt genau  darstellen, dann Repärsentationen / Ausgabeformaten / Präsentationsformen in HTML, als PDF etc. automatisch generieren (sog. \emph{single source} Prinzip).

Vom Buch zu unseren digitalen Daten, von den digitalen Daten wieder zum Buch.
Bei der Konvertierung geht Information eventuell verloren. Z.B. "in Anführungsstrichen", was bedeutet das jetzt? Der Werktitel? Was, wenn er dann kursiv genmacht werden soll? Wenn kein Markup mehr dabei ist, wird es schwerer, das automatisiert zu ändern. Daher auch möglichst keine Darstellunganweisenen (`in quotes') hardcoden (= direkt reinschreiben), sondern lieber in den Basisdaten, z.B. als \texttt{<q>} (\emph{quote}).
\end{frame}

\begin{frame}[allowframebreaks]{Datenmodellierung}

\bg{alert}{white}{Modellbegriff}
„Alle Erkenntnis ist Erkenntnis in Modellen oder durch Modelle und jede
menschliche Weltbegegnung überhaupt bedarf des Mediums Modell“ (Herbert Stachowiak 1973)
\smallskip

\bg{alert}{white}{Abbildungsmerkmal}
Ein Modell ist immer ein Abbild von etwas, eine Repräsentation
natürlicher oder künstlicher Originale.
\smallskip

\bg{alert}{white}{Verkürzungsmerkmal}
Ein Modell erfasst nicht alle Attribute des Originals, sondern nur
diejenigen, die dem Modellschaffenden bzw. der Nutzerin/dem Nutzer
relevant erscheinen. Ein Modell abstrahiert. Z.B. kann man eine Statue immer nur von einer Seite fotographieren (bzw. sogar sehen! D.h. ich sehe nie `das Ding an sich', sondern einen Ausschnitt aus der Summe aller Eigenschaften, die es ausmachen).
\smallskip

\bg{alert}{white}{Pragmatisches Merkmal}
Orientierung am Nützlichen. Ein Modell wird vom Modellschaffenden
bzw. der Nutzerin/dem Nutzer innerhalb einer bestimmten Zeitspanne
und zu einem bestimmten Verwendungszweck an Stelle eines Originals
eingesetzt.
\smallskip

\bg{alert}{white}{Datenmodellierung}
Modell =  Ausschnitt aus der realen Welt, allerdings
werden im Modell nur jene Attribute berücksichtigt, die für
meine Fragestellung relevant sind.
 Modell und Realwelt weichen somit voneinander ab.
\smallskip

\bg{alert}{white}{Modell = Abstraktion (=Klasse)}
Konkretisierungen (Kochrezepte, Briefe, Gedichte usw.) eines
Modells werden als Instanzen bezeichnet.
\smallskip

Standardisierte Modelle ermöglichen die gemeinsame
Auswertung, datenübergreifende Suche, Datenaustausch\dots
\smallskip

\bg{alert}{white}{Modelle = vereinfachte Repräsentationen von Teilen der realen Welt.}

\end{frame}

\begin{frame}[allowframebreaks]{McCarty}

\bg{alert}{white}{Willard McCarty, „Modeling: A Study in words and Meanings“ (2004)}
\begin{enumerate}
    \item Ein Modell ist eine Repräsentation von etwas, z.B. für Studienzwecke.
    \item Modellierung = heuristischer Prozess in dem Modelle kreiert und verwendet werden, um Probleme zu lösen.
    \item \emph{model of} (beschreibend) vs. \emph{model for} (z.B. Plan zum Hausbau).
    \item Modellierung als experiementeller Vorgang mit dem Ziel der Welterfassung.
\end{enumerate}
\framebreak

\bg{w3schools}{white}{Modellierung als iterativer Prozess} \\
\begin{itemize}
    \item Bestehendes Wissen einbringen
    \item Modell bilden
    \item Das Modell verwenden um das Thema zu untersuchen
    \item an den Diskrepanzen zw. Modell und Realität zeigt sich der Verfeinerungsbedarf, da das Modell ja immer nur eine Annährung ist. Diese zeigt aber wieder, wo noch nachgebessert werden muss. 
    \item mit neuem Wissen aus Modell neu modellieren
\end{itemize}

\end{frame}





\input{xpath}




\section{Transformationen mit XSLT}

\begin{frame}[fragile,allowframebreaks]{Was ist XSLT?}
\bg{alert}{white}{XSL} (\emph{eXtensible Stylesheet Language}) ist eine Programmiersprache zur Transformation von XML-Daten. Diese erlaubt die Speicherung von abstrakten Basisdaten, aus denen dann unterschiedlichste Repräsentationen, z.B. die Strukturinformationen für HTML, automatisch generiert werden können (\emph{single source}-Prinzip).  \textbf{XSLT} (=XSL-Transformations) wird häufig synonym dazu verwendet, ist aber eigentlich nur ein Teil neben \textbf{XSL-FO} (\emph{Formatting Objects}, zur Beschreibung von Druckdokumenten als PDF), das seit 2012 nicht mehr weiterentwickelt wird.

\bg{alert}{white}{XSLT} besteht aus Verarbeitungsmustern (sog. `Schablonen' oder \emph{templates}), wobei das XML-Dokument durchgegangen wird und Schablonen angewendet werden, sobald eine auf den aktuell angefundenen Inhalt passt. Die Verarbeitung durch den Parser beginnt beim Wurzelknoten, weswegen es immer eine erste Schablone gibt, die auf alles passt:
\begin{verbatim}
    <xsl:template match="/">...</xsl:template>
\end{verbatim}
Wird für angetroffene Inhalte kein Template gefunden, so werden die \emph{Built-in}-Regeln angewendet, die die Inhalte der Elemente (inkl. Unterelemente!) ausgeben. Dies begegnet einem oft in dem Umstand, dass einfach der ganze Dokumenteninhalt drauflosgeprinted wird, was oft nicht eigentlich erwünscht ist. Z.B. möchte man oft ja nur den \texttt{body}-Inhalt eines Dokuments und nicht alle Metainformationen des \texttt{teiHeader}s unsortiert ausgedruckt haben. Dieses grundlegende Verhalten muss man daher bewusst unterbinden. 

\end{frame}

%-----------------------------------------
\begin{frame}[allowframebreaks]{XSLT-Intro}

\bg{alert}{white}{XSL(T): eXtensible Stylesheet Language (Transformations)}\\
\bgupper{w3schools}{black}{.xsl}\\

Ein XSLT-Stylesheet (ist selbst ein XML-Dokument und) beschreibt Regeln für den  Transformationsprozess einer Eingabedatei (XML-Dokument) in eine oder
mehrere Ausgabedateien (XML, XHTML, Text). 
\smallskip

\bg{alert}{white}{Ausgabeformate sind z.B.\dots}
\begin{itemize}
    \item \textbf{(x)HTML} \item \textbf{XML} $\to$ z.B. Word, TEI und andere XML-Standards, RDF, SVG \item \textbf{Text} $\to$ z.B. LaTeX ($\to$ PDF), RTF, MARC.
\end{itemize}

\alert{Warnung am Rande:} \LaTeX{} z.B. ist auch Markup, hat aber andere geschützte Entitäten, bzw. andere Konventionen, wie man Dinge, wie etwa den Ampersand escaped. \texttt{xsl:output} und nachkontrollieren!

\end{frame}

\begin{frame}[allowframebreaks]{XSLT Ressourcen}

\href{https://www.i-d-e.de/?s=referenz}{IDE Kurzreferenzen}
\sep 
\href{https://www.w3schools.com/xml/xsl_intro.asp}{W3Schools-Tutorial}

\end{frame}



\begin{frame}[fragile,allowframebreaks]{Einstieg}
In Oxygen sowohl ein zu transformierendes XML als auch ein neues XSL-Stylesheet eröffnen.

\begin{myxml}{Auto-generiertes Start-Dokument}
<?xml version="1.0" encoding="UTF-8"?>
<xsl:stylesheet xmlns:xsl="http://www.w3.org/1999/XSL/Transform"
    xmlns:xs="http://www.w3.org/2001/XMLSchema"
    exclude-result-prefixes="xs"
    version="2.0">
    
</xsl:stylesheet>
\end{myxml}
\framebreak

Dann in die XML-Datei reinklicken und ein Transformationsszenario konfigurieren (Schraubenzieher-Symbol):
\begin{enumerate}\footnotesize
    \item `Neu' auswählen (XSLT-Transformation).
    \item Reiter XSLT: XML-Datei aussuchen, XSL-Datei aussuchen (kann man sonst auch auf massenweise Dokumente anwenden (Projekt-Optionen).
    \item Transformator wählen: Saxon (9er Version, sonst egal welcher).
    \item Reiter XSL-FO: Hier egal, weil das eine andere Transformation als XSLT ist (zum PDFs machen, was alternativ zu \LaTeX{} dort gehen würde).
    \item Reiter Ausgabedatei: `Datei speichern unter' $\to$ Grünen Pfeil anklicken (\verb|${cfn}| auswählen. Dann \texttt{-transform.xml} hinzufügen, damit die Orignaldatei nicht überschrieben wird. (Bei erneuter Transformation wird diese Datei immer wieder überschrieben, außer man benennt sie um).
    \item Wahlweise `Im Editor öffnen' und anzeigen als XML, falls XML oder XHTML, falls HTML.
    \item Ok, dann anwenden.
\end{enumerate}

Damit wird nun unser, bishe noch sehr leeres, Stylesheet auf die Datei angewendet. \alert{ Was fällt auf? }

Das sich zeigende Verhalten ist das Standard-Verhalten von XSLT, wenn es keine zutreffenden Verarbeitungsregeln findet.
\end{frame}


\begin{frame}[fragile,allowframebreaks]{Standard-Transformationsablauf}
\begin{myxml}{Leeres Match-Root-Template}
<?xml version="1.0" encoding="UTF-8"?>
<xsl:stylesheet xmlns:xsl="http://www.w3.org/1999/XSL/Transform"
    xmlns:xs="http://www.w3.org/2001/XMLSchema"
    exclude-result-prefixes="xs"
    version="2.0">
       
    <xsl:template match="/">
        <!-- und jetzt? -->
    </xsl:template>
    
</xsl:stylesheet>
\end{myxml}
Mit dem folgenden äußern wir, dass wir möchten, dass die Standard-Verarbeitungsweise wieder verwendet wird. Der angefundene Text wird ohne Elemente hineingeschrieben. Achtung, das ist jetzt nicht eigentlich ein XML-Dokument, weil es keine Elemente oder hierarchische Struktur hat!

\end{frame}

\begin{frame}[fragile,allowframebreaks]{Standard-Transformationsablauf}
\begin{myxml}{Standard-Regeln mit \texttt{apply-templates}}

<?xml version="1.0" encoding="UTF-8"?>
<xsl:stylesheet xmlns:xsl="http://www.w3.org/1999/XSL/Transform"
    xmlns:xs="http://www.w3.org/2001/XMLSchema"
    exclude-result-prefixes="xs"
    version="2.0">
       
    <xsl:template match="/">
        <xsl:apply-templates/>
    </xsl:template>
    
</xsl:stylesheet>
\end{myxml}

\begin{myxml}{Struktur `hardcoden' = händisch reinschreiben}
<?xml version="1.0" encoding="UTF-8"?>
<xsl:stylesheet xmlns:xsl="http://www.w3.org/1999/XSL/Transform"
    xmlns:xs="http://www.w3.org/2001/XMLSchema"
    exclude-result-prefixes="xs"
    version="2.0">
       
    <xsl:template match="/">
        <xml>
            <xsl:apply-templates/> 
        </xml>
    </xsl:template>
    
</xsl:stylesheet>
\end{myxml}
\end{frame}

\begin{frame}[fragile,allowframebreaks]{XSLT-Standardablauf II}
Statt die Elementenhierarchie mühsam händisch wieder reinschreiben zu müssen, können wir auch Templates (`Schablonen') definieren, die immer aktiv werden, wenn ein bestimmtes Element gefunden wird. Lies: Wann immer Du \texttt{<p/>} findest, mach XY.
Wir deklarieren vorher noch einen Namespace (\texttt{xmlns}), damit unsere TEI-Dateien erkannt werden:
\begin{verbatim}
    xmlns:t="http://www.tei-c.org/ns/1.0"
\end{verbatim}
Elemente müssen jetzt immer mit \texttt{t:} davor angesprochen. \texttt{t:p} lies: \texttt{<p>} aus dem Namespace, der mit \texttt{t:} abgekürzt ist (siehe Stylesheet-Deklaration oben).
\begin{myxml}{\texttt{p}-Elemente matchen}
<?xml version="1.0" encoding="UTF-8"?>
<xsl:stylesheet xmlns:xsl="http://www.w3.org/1999/XSL/Transform"
    xmlns:xs="http://www.w3.org/2001/XMLSchema"
    xmlns:t="http://www.tei-c.org/ns/1.0"
    exclude-result-prefixes="xs"
    version="2.0">
       
    <xsl:template match="/">
        <xml>
            <xsl:apply-templates/> 
        </xml>
    </xsl:template>
    
    <xsl:template match="t:p">
        <neu>
            <xsl:apply-templates/>
        </neu>
    </xsl:template>
    
</xsl:stylesheet>
\end{myxml}

\end{frame}

\begin{frame}[fragile,allowframebreaks]{XSLT-Basics III}
Alle \texttt{p}s, inklusive ihrem Inhalt (dank \texttt{apply-templates}) werden erhalten, aber sie heißen nun statt \texttt{<p> <neu>} (sinnfrei \& zu Demonstrationszwecken).
Dieses Vorgehen jedenfalls nennt sich das sogenannte \bg{alert}{white}{Push-Paradigma}.


Das Gegenstück dazu, das \bg{alert}{white}{Pull-Paradigma} sehen wir jetzt beim hinzugefügten \texttt{<head>}. In dieses neu erzeugte Element holen wir uns bewusst den Inhalt vom \texttt{<title>}-Element rein (Achtung, hier XPath-Wissen anwenden: Wenn es mehrere \texttt{<title>}-Elemente gibt, nimmt er \emph{alle}. Im Zweifelsfall explizit und mist absoluten Pfaden arbeiten. Außerdem ist \texttt{value-of select=""} nicht dasselbe wie \texttt{apply-templates}. \texttt{apply-templates} schaut erst, ob es (weiter \emph{unten} im Code!) noch Regeln gibt, die auf die Elemente zutreffen, die im gefundenen Ding eingeschachtelt sind und wendet ggf. diese an. \texttt{value-of select=""} kopiert nur den aktuell angesprochenen Textinhalt; tiefere Verschachtelungen gehen verloren. \texttt{value-of select="."} gibt den aktuellen Knoteninhalt (Text). \verb|attribut={@rend}| ist eine Kurzschreibweise dafür bei Attributen.
\begin{myxml}{}
<?xml version="1.0" encoding="UTF-8"?>
<xsl:stylesheet xmlns:xsl="http://www.w3.org/1999/XSL/Transform"
    xmlns:xs="http://www.w3.org/2001/XMLSchema"
    xmlns:t="http://www.tei-c.org/ns/1.0"
    exclude-result-prefixes="xs"
    version="2.0">
       
    <xsl:template match="/">
        <xml>
            <head><xsl:value-of select="//t:title"/></head>
            <xsl:apply-templates/> 
        </xml>
    </xsl:template>
    
    <xsl:template match="t:p">
        <neu attribut={@rend}>
            <xsl:apply-templates/>
        </neu>
    </xsl:template>
    
</xsl:stylesheet>
\end{myxml}

\end{frame}

\begin{frame}[fragile,allowframebreaks]{Inhalte löschen}
Per Automatismus werden in XSLT alle Inhalte (Werte) erhalten -- aber wir haben schon gesehen, dass man mit leeren Regeln bewusst löschen kann. Wir löschen im Beispiel also probehalber alle \texttt{<hi>} (inklusive Inhalt!). Hier darauf achten, dass man nicht versehentlich auf die Quelldatei schreibt, sonst ist alles weg (siehe Transformationsszenario konfigurieren).

\begin{myxml}{Bewusstes Löschen durch leere Regel}
    <xsl:template match="t:hi">
        <!-- loeschen -->
    </xsl:template>
\end{myxml}
\framebreak

Das sind im Grunde schon alle Regeln, die man wissen muss. Daneben gibt es natürlich noch mehr Funktionen, falls man kompliziertere Sachen machen will. Diese finden sich im Folgenden erklärt, werden aber evtl. gar nicht gebraucht.
\medskip

Folgend noch ein Beispiel\dots 

\end{frame}

\begin{frame}[fragile,allowframebreaks]{Beispiel}
\begin{myxml}{Gedicht $\to$ HTML-Webansicht}
<?xml version="1.0" encoding="UTF-8"?>
<xsl:stylesheet xmlns:xsl="http://www.w3.org/1999/XSL/Transform"
    xmlns="http://www.w3.org/1999/xhtml" version="2.0">
    
    <xsl:template match="/">
        <html>
            <head>
                <title><xsl:value-of select="div/head" /></title>
            </head>
            <body>
                <xsl:apply-templates select="div" />
            </body>
        </html>
    </xsl:template>
    
    <xsl:template match="div">
        <h1><xsl:value-of select="head" /></h1>
        <div>
            <xsl:apply-templates select="lg" />
        </div>
    </xsl:template>
    
    <xsl:template match="lg">
        <p><xsl:apply-templates /></p>    
    </xsl:template>
    
    <xsl:template match="l">
        <xsl:value-of select="." />
        <br />
    </xsl:template>
</xsl:stylesheet>
\end{myxml}
\end{frame}

\begin{frame}[fragile,allowframebreaks]{Weitere Funktionen}
\bg{alert}{white}{Variablen}
\begin{verbatim}
    <xsl:variable name="platzhalter" select="inhalt">
\end{verbatim}
XSLT führt bei Variablen einen sog. \emph{pass by value}, nicht \emph{by reference} aus, d.h. es wird nur der Wert übergeben, den das ausgewählte Element zu einem Zeitpunkt hat. Wenn ich die Variable zu früh setze, wähle ich dabei womöglich einen fixierten Wert. Also am besten immer so nah wie möglich am `Ziel'. 
\smallskip

\framebreak

\bg{alert}{white}{Attribute}
\texttt{<xsl:attribute>} dient zum Setzen von Attributen im Resultat-Dokument. Dies ist nicht nötig, wenn man nur einen Inhalt will, dann kann auch die Kurzschreibweise verwendet werden.
\smallskip

\bg{alert}{white}{\emph{plain text}}
\texttt{<xsl:text>} setzt Text `as is', d.h. fügt nicht unerwünschte Leerzeichen ein, was sonst passieren. Überhaupt geht XSLT mit Whitespace relativ frei um, außer man verbietet es explizt (\texttt{preserve-space}.
\smallskip

\bg{alert}{white}{Sortieren}
Es gibt auch \texttt{<xsl:sort>} zum Sortieren.
\smallskip

\bg{alert}{white}{Bedingungen}
\texttt{<xsl:choose>} zum Auswählen unterschiedlicher Verarbeitung je nachdem, was angefunden wird.
\texttt{<xsl:if>} führt Anweisungen nur aus, wenn eine Bedingung erfüllt ist (z.B. falls es Element XY gibt, oder ein Element Inhalt XY hat).


\framebreak

\begin{myxml}{Personenliste in HTML mit \texttt{for-each}}
<ul> <!-- ungeordnete Liste in HTML -->
<xsl:for-each select="t:persName">
    <li> <!-- Listenelement -->
    <xsl:value-of select="." />
    </li>
</xsl:for-each>
</ul>
\end{myxml}

\begin{myxml}{Attribute setzen}
<xsl:for-each select="t:persName">
    <span class={@rend}><xsl:value-of select="."></span>
</xsl:for-each>
<!-- ODER -->

<span>
    <xsl:attribute name="id">
        <xsl:value-of select="@rend"><xsl:text>-person</xsl:text>
    <xsl:attribute>
    <xsl:attribute name="interpretation">
        <xsl:value-of select="concat(@rend,@ana)">
    <xsl:attribute>
    <xsl:apply-templates /> <!-- für den eig. Inhalt -->
<span>
<!-- Resultat z.B. 
    <span id="monster-person" interpretation="monster-boese">
    Weird Sisters </span> -->
\end{myxml}

\begin{myxml}{For-each-group}

<!-- könnte man auch nach Alphabet sortieren oder so. Nur gezielt verwenden, macht in komplexen Dokumenten oft ominöse Fehler. -->

<xsl:for-each-group select"dings" group-starting-with="dings[@rend="Startdings]]"
    <xsl:apply-templates select="current-group()">
    <xsl:value-of select="current()/dingens">
</xsl:for-each-group>

\end{myxml}

\begin{myxml}{Mehrere Quelldokumente zusammenführen (\texttt{document()})}
<xsl:apply-templates select="document('Letter1_TEI.xml')/tei:TEI/tei:text/tei:body/tei:div"/>
\end{myxml}
\end{frame}

\begin{frame}[fragile,allowframebreaks]{\emph{flow control} / Bedingungen}


Beispiel: Wähle jede Person (\texttt{//persName}) und generiere eine Aufzählung (\texttt{<ul>}) der Nachnamen (\texttt{lastname}):
\begin{myxml}{Schleifen: for-each}
<ul>
<xsl:for-each select="//persName">
    <xsl:sort select="lastname" order="ascending" />
    <li> <xsl:value-of select="lastname"/> </li>
</xsl:for-each>
</ul>
\end{myxml}


\texttt{xsl:if} führt Anweisungen nur aus, wenn die Bedingung in \texttt{@test} erfüllt ist.
\begin{myxml}{Bedingungen I: Falls}
<xsl:if test=" xpath-ausdruck "> ... </xsl:if>

<xsl:for-each select="//book">
    <xsl:if test=" author = 'Cicero' ">
        <li><xsl:value-of select="title"/></li>
    </xsl:if>
</xsl:for-each>
\end{myxml}

Auch die Unterscheidung mehrerer Fälle ist möglich. So kann ich z.B. für das HTML unterschiedlichen Personen unterschiedliche Farben zuweisen. Dazu frage ich, in \texttt{@test} welche ID ein angetroffener \texttt{persName} hat und verarbeite dann je nachdem unterschiedlich.
\begin{myxml}{Bedingungen II: Auswählen}
<xsl:choose>
    <xsl:when test=" xpath-ausdruck "> ... </xsl:when>
    <xsl:otherwise> ... </xsl:otherwise>
</xsl:choose>
\end{myxml}

Man kann zudem sortieren (\texttt{xsl:sort}), kopieren (\texttt{xsl:copy - xsl:copy-of})und Variablen verwenden. 

\end{frame}

\begin{frame}[fragile,allowframebreaks]{XSL-Paradigmen: Push}
\begin{myxml}{Push-Paradigma}
<xsl:apply-templates select="xpath-ausdruck"/>
\end{myxml}
Geeignet vor allem: Für das Verarbeiten des Text-Bodys. Hier werden angetroffene Dinge einfach verarbeitet, wann immer sie auftreten. Man muss deren Reihenfolge nicht kennen. Die Dokumentenstruktur des `Herkunftsdokuments' wird weitgehend erhalten (soweit es eben in den Template-Regeln definiert ist).

\begin{myxml}{Push-Paradigma}
<xsl:template match="/">
    <xsl:apply-templates/>
</xsl:template>
\end{myxml}

\bg{alert}{white}{Push Processing}
Jedes Element bekommt eine Regel oder \emph{built-ins} werden verwendet. Nützlich falls Input- und Outputdokument dieselbe Struktur haben sollen. 


\begin{myxml}{Call Template / Push-Methode}
<xsl:template match="/">
    <xsl:call-template name="irgendein-name ">
</xsl:template>
<xsl:template name="irgendein-name">
    ... macht etwas ...
</xsl:template>
\end{myxml}

\end{frame}


\begin{frame}[fragile,allowframebreaks]{XSL-Paradigmen: Pull}

\begin{myxml}{Pull-Paradigma}
<xsl:call-templates name="name-eines-templates "/>
\end{myxml}
Geeignet vor allem: Herausziehen von Metadaten aus dem \texttt{teiHeader}, um sie irgendwo extra anzuzeigen.
Nochmals `über das Dokument drübergehen' und z.B. einen Personenindex machen, den man an einer bestimmten Stelle hingesetzt haben will (\texttt{<xsl:for-each>} ist z.B. ein typischer Befehl der Push-Methode.

\framebreak

\bg{alert}{white}{Pull Processing}
Gewisse Knoten explizit auswählen. Kontrolle darüber, nur bestimmte Knoten im Output-Dokument weiterzuverwenden.
Gut geeignet, falls das Output-Dokument eine völlig andere Struktur haben soll, als der Inhalt oder nur sehr selektiv Inhalte übernommen werden sollen (z.B. nur eine Liste aller enthaltenen Personennamen braucht nicht den ganzen restlichen Text dazu verarbeiten).

\begin{myxml}{Pull-Paradigma}
<xsl:value-of select=’’pattern’’/>
<xsl:apply-templates select=’’pattern’’/>
<xsl:for-each select=’’pattern’’/>
\end{myxml}

\end{frame}

\begin{frame}[fragile,allowframebreaks]{Paradigmen in der Praxis}
In der Praxis wird meist eine Kombination von beiden verwendet: Die Metadaten des \texttt{teiHeader} werden, z.B. in der Form eines Zitiervorschlags, mit dem Pull-Paradigma ausgelesen. Dann geht man an die Stelle des Output-Dokuments, wo der hauptsächliche Inhalt (\emph{body}) hin soll und wählt aus, dass gar nicht mehr das ganze Dokument automatisch per Pull-Paradigma verarbeitet werden soll, sondern nur alles ab dem \texttt{body}.
\begin{myxml}{apply-templates select}
<xsl:apply templates select="//t:body" />
<xsl:apply-templates select="head" mode="toc"/>
\end{myxml}

\alert{Anderes Beispiel:} Kapitelüberschriften innerhalb des Buches per Push-Paradigma, für das Inhaltsverzeichnis im Pull-Paradigma. Mithilfe von \texttt{mode}s kann man zudem noch unterscheiden, dass in gewissen Situationen andere Templateregeln angewandt werden. 

\verb|<xsl:template match="">| verarbeitet ein Element, wo immer es angetroffen wird.

\verb|<xsl:value-of select="">| gibt den Inhalt eines Elements als String aus -- wenn noch Kindelemente sind, gehen die Knoten/Auszeichnungen verloren, weil nur mehr plain text kopiert wird.

\texttt{<xsl:apply-templates>} geht die Regeln durch, wo er dann alles findet, was unter \texttt{"match"} definiert ist  -- er sucht aber Unter-Verschachtelungen nur weiter, wenn man das explizit so angibt (mit weiterem \texttt{apply-templates}).

\end{frame}

\begin{frame}[fragile,allowframebreaks]{Facts}
%\bg{alert}{white}{XSLT als Sprache} ist deklarativ (wie CSS nur mächtiger) und dynamically-typed, d.h. es arbeitet eher mit Werten als mit Variablen (wie JavaScript).


\bg{alert}{white}{Vorgehen bei Stylesheets}
Regelwerk nach Top-down-Prinzip, ausgehend vom
Wurzelelement des Quelldokuments.

\bg{alert}{white}{Wurzelelement \texttt{xsl:stylesheet}}
XSLT ist selbst XML, daher hat es auch ein alles umgebendes Wurzelelement. Hier werden einige Dinge deklariert (Namensräume, XSLT-Processor, etc.).

\end{frame}



 
\end{document}
